% ============================================================================
% Fichier  : partie2/P02_C07_mission_principes_cycle.tex
% Titre    : La Mission d'Audit : Principes et Cycle
% Auteur   : MINKA MI NGUIDJOÏ Thierry Emmanuel
% Version  : 1.0 -- novembre 2026
% ============================================================================
% Base : ACSI C05 (sept principes ISO 19011:2018 ; charte IIA 2024 ; cycle
% en 5 phases ; outils O0-O9 ; calendrier AUDIX CONSEIL 4 mai - 17 juillet
% 2026 ; domaine B 6/2/4 puis 8/3/1).
% Sources : ISO 19011:2018 ; Normes mondiales d'audit interne de l'IIA
% (2024) ; ITAF (ISACA, 2020) pour la supervision et la documentation ;
% articles A0 (axiomes 5 et 6) et A1.
% Corrigés : corriges/Corr_C07.tex.
% ============================================================================

\chapter{La Mission d'Audit : Principes et Cycle}
\label{chap:mission}
\chaptermark{La mission d'audit}

\epigraph{%
    « Une mission d'audit se gagne avant le terrain\\
    et se perd après le rapport. »%
}{Adage des chefs de mission}

\niveauChapitre{Fondamental}{%
    Maîtrise du Chapitre~\ref{chap:audit-auditologie} (audit, assurance,
    axiomes de l'auditologie)%
}

\begin{boxobjectifs}
\begin{itemize}
    \item Énoncer les sept principes de l'audit selon l'ISO~19011:2018 et
          les illustrer.
    \item Rédiger les rubriques d'une charte d'audit et d'une lettre de
          mission.
    \item Répartir les rôles dans une équipe d'audit et organiser la
          supervision.
    \item Décrire les cinq phases de la mission, leurs entrées et leurs
          livrables, et y situer les dix outils O0 à O9.
    \item Tenir des papiers de travail référencés et revus.
    \item Organiser le contradictoire, le relier à l'opposabilité des
          conclusions (axiome~6) et mesurer ce qu'il apporte à la
          connaissance.
\end{itemize}
\end{boxobjectifs}

\bigskip

% ============================================================================
\section*{Question directrice : comment ne pas refaire la première mission~?}
% ============================================================================

La première mission de LiZbiZ (Chapitre~\ref{chap:prologue}) a échoué
faute de mandat, de critères et de preuves. Pour la seconde, le cabinet
AUDIX CONSEIL n'improvise rien~: M.~BEYALA ouvre la mission par une réunion
avec le directeur général, remet une lettre de mission, annonce un
calendrier en cinq phases et prévient que chaque constat sera discuté avec
les audités avant d'être écrit.

\medskip
\textbf{Plan du chapitre.}
\begin{enumerate}
    \item Les principes de l'audit (section~\ref{sec:c07-principes}).
    \item La charte et la lettre de mission (section~\ref{sec:c07-charte}).
    \item L'équipe et ses responsabilités (section~\ref{sec:c07-equipe}).
    \item Le cycle de la mission (section~\ref{sec:c07-cycle}).
    \item La boîte à outils (section~\ref{sec:c07-outils}).
    \item Les papiers de travail (section~\ref{sec:c07-papiers}).
    \item Le contradictoire (section~\ref{sec:c07-contradictoire}).
\end{enumerate}
Les Chapitres~\ref{chap:planifier} à~\ref{chap:suivi} détaillent ensuite
chaque phase.

\separateur

% ============================================================================
\section{Les Principes de l'Audit}
\label{sec:c07-principes}
% ============================================================================

L'ISO~19011 donne des lignes directrices pour l'audit des systèmes de
management. Ses principes valent pour tout audit, et en particulier pour
l'audit des SI~\cite{iso19011_2018}.

\begin{tableaucours}{L{3.4cm}L{4.4cm}L{4.4cm}}{%
    \tableauHeader{Principe} &
    \tableauHeader{Sens (reformulé)} &
    \tableauHeader{Chez LiZbiZ}
}
Intégrité & Honnêteté, diligence, responsabilité & Ne pas taire un constat
qui gêne la direction \\
Présentation impartiale & Rendre compte avec exactitude et fidélité &
Ne pas grossir un risque mineur, ne pas minimiser un risque majeur \\
Conscience professionnelle & Diligence et jugement proportionnés à
l'enjeu & Ne pas lancer un test de charge à l'heure de pointe \\
Confidentialité & Protéger l'information obtenue & Ne pas divulguer les
failles constatées \\
Indépendance & Impartialité, absence de conflit d'intérêts & Ne pas auditer
l'application qu'on a conçue \\
Approche fondée sur la preuve & Des conclusions vérifiables, appuyées sur
un échantillon adapté & Une capture, une extraction, pas un ouï-dire \\
Approche par les risques & Concentrer l'effort sur ce qui compte pour les
objectifs & Commencer par les risques intolérables du registre \\
\end{tableaucours}
\vspace{-0.8em}
\begin{center}{\small\itshape Tableau~7.1. Les sept principes de l'audit
selon l'ISO~19011:2018 (notre reformulation)}\end{center}

\begin{boxinfo}
\textbf{Six ou sept principes~?} L'édition de 2011 de l'ISO~19011 en
comptait six. L'édition de 2018, en vigueur, ajoute l'\emph{approche par
les risques}. Un support qui annonce « six principes » se réfère à
l'ancienne édition.
\end{boxinfo}

Deux de ces principes trouvent un fondement dans l'auditologie. L'approche
fondée sur la preuve découle de l'axiome d'opposabilité~: une conclusion
doit pouvoir être défendue devant un tiers. L'indépendance, on l'a vu au
Chapitre~\ref{chap:audit-auditologie}, se déduit de l'interaction
stratégique et de l'opposabilité~\cite{minka2026auditologie}.

\subsection{Les principes en conflit}

Les principes ne se concilient pas toujours d'eux-mêmes. La
confidentialité interdit de divulguer une faille~; la présentation
impartiale oblige à la signaler. La solution est dans les destinataires~:
la faille est décrite précisément au commanditaire, et de façon générale
dans toute version du rapport qui circule plus largement. L'approche par les
risques invite à négliger les domaines peu risqués~; la conscience
professionnelle interdit de les ignorer s'ils présentent un signal. La
règle du cours est d'écrire, dans le plan d'audit, ce qui est exclu et
pourquoi.

% ============================================================================
\section{Charte d'Audit et Lettre de Mission}
\label{sec:c07-charte}
% ============================================================================

\subsection{La charte d'audit}

\begin{boxdef}
\termecle{Charte d'audit}\index{charte d'audit}~: document approuvé par
l'organe de gouvernance, qui fixe la mission de l'audit, ses pouvoirs (accès
aux personnes, aux documents et aux systèmes), son rattachement, son
indépendance et ses responsabilités.
\end{boxdef}

Les Normes mondiales d'audit interne de l'IIA placent sous la
responsabilité du conseil la supervision de la fonction d'audit, son mandat
et les moyens de son indépendance~\cite{iia2024gias}. La charte en est
l'instrument~: sans elle, l'auditeur dépend du bon vouloir de ceux qu'il
audite.

\begin{tableaucours}{L{3.2cm}L{9.0cm}}{%
    \tableauHeader{Rubrique} &
    \tableauHeader{Contenu attendu}
}
Mission & Finalité de l'audit (assurance et conseil), périmètre de
l'organisation couvert \\
Rattachement & Rattachement fonctionnel au comité d'audit ou au conseil,
rattachement administratif à la direction générale \\
Pouvoirs & Accès sans restriction aux personnes, aux documents, aux biens et
aux systèmes, y compris en lecture sur les bases de production \\
Indépendance & Absence de responsabilité opérationnelle, déclaration
annuelle d'indépendance, règles en cas de conflit d'intérêts \\
Responsabilités & Plan annuel fondé sur les risques, rapports, suivi des
recommandations, normes appliquées \\
Relations & Avec le commissaire aux comptes, les régulateurs, les
prestataires d'audit \\
\end{tableaucours}
\vspace{-0.8em}
\begin{center}{\small\itshape Tableau~7.2. Les rubriques d'une charte
d'audit}\end{center}

\subsection{La lettre de mission}

Chaque mission est ensuite fondée sur une \termecle{lettre de
mission}\index{lettre de mission}, ou des termes de référence, qui précise
ce qu'elle doit démontrer. Elle répond à quatre questions~: quoi (objectifs
et périmètre), selon quoi (référentiels et critères), comment (approche,
calendrier, équipe) et pour qui (destinataires et livrables). Le
Chapitre~\ref{chap:audit-auditologie} en a donné un exemple pour LiZbiZ~;
le Chapitre~\ref{chap:planifier} en détaille le contenu.

\begin{boxalert}
\textbf{Violation typique.} Le plan d'audit interne d'une entreprise est
approuvé par le directeur général, dont relèvent les directions auditées.

\smallskip
\textbf{Règle~:} le plan d'audit et les moyens de l'audit sont approuvés par
l'organe de gouvernance (conseil, comité d'audit), et non par la direction
auditée.
\end{boxalert}

% ============================================================================
\section{L'Équipe et ses Responsabilités}
\label{sec:c07-equipe}
% ============================================================================

\subsection{Quatre rôles}

\begin{tableaucours}{L{2.8cm}L{5.0cm}L{4.4cm}}{%
    \tableauHeader{Rôle} &
    \tableauHeader{Responsabilités} &
    \tableauHeader{Chez AUDIX CONSEIL}
}
Associé signataire & Accepte la mission, signe le rapport, répond de la
qualité devant le client & L'associée du cabinet, Mme~TCHOUA \\
Chef de mission & Planifie, répartit les travaux, revoit les papiers de
travail, conduit le contradictoire & M.~BEYALA \\
Auditeur & Exécute les tests, tient les papiers de travail, rédige les
projets d'observation & Deux auditeurs seniors, un junior \\
Expert & Apporte une compétence technique (réseau, base de données),
sous la responsabilité du chef de mission & Un ingénieur réseau pour les
centres de données \\
\end{tableaucours}
\vspace{-0.8em}
\begin{center}{\small\itshape Tableau~7.3. Les rôles dans une équipe
d'audit}\end{center}

\subsection{La supervision}

Chaque papier de travail est revu par une personne plus expérimentée que
son auteur. La revue porte sur trois points~: le test répond-il au
critère~? la preuve suffit-elle à la cote~? la conclusion est-elle séparée
de la constatation~? La revue laisse une trace (nom, date, remarques et leur
traitement). Le cadre professionnel de l'ISACA fait de cette supervision
une exigence~\cite{isaca2020itaf}.

\begin{boxalert}
\textbf{Violation typique.} Un expert externe remet un rapport technique que
l'équipe recopie dans le rapport d'audit sans l'avoir compris.

\smallskip
\textbf{Règle~:} le chef de mission répond des travaux de l'expert comme
des siens. Il en vérifie la méthode, la preuve et la cohérence avec les
critères de la mission avant de les utiliser.
\end{boxalert}

% ============================================================================
\section{Le Cycle de la Mission}
\label{sec:c07-cycle}
% ============================================================================

\subsection{Cinq phases}

\noindent\begin{minipage}{\linewidth}
\begin{center}
\begin{tikzpicture}[
    phase/.style = {
        draw=CouleurPrimaire, fill=CouleurPrimaire, rounded corners=4pt,
        text width=2.1cm, minimum height=1.1cm, align=center,
        font=\scriptsize\bfseries\color{white}
    },
    livr/.style = {
        draw=CouleurBordDef, fill=CouleurDef, rounded corners=3pt,
        text width=2.0cm, align=center, font=\scriptsize\color{CouleurPrimaire}
    },
    fleche/.style = {-{Stealth[length=5pt]}, thick, color=CouleurSecondaire}
]
\foreach \n/\t/\l [count=\i] in {
    0/Préparation/Termes de référence,
    1/Planification/Plan d'audit,
    2/Exécution/Papiers de travail,
    3/Rapport/Rapport contradictoire,
    4/Suivi/Tableau de suivi} {
    \node[phase] (p\i) at (2.55*\i, 0) {Phase~\n\\ \t};
    \node[livr, below=0.35cm of p\i] {\l};
}
\foreach \i [evaluate=\i as \j using int(\i+1)] in {1,...,4}
    \draw[fleche] (p\i) -- (p\j);
\draw[fleche] (p5.north) to[out=150, in=30] node[above, font=\scriptsize\itshape]
    {mission suivante} (p1.north);
\end{tikzpicture}
\end{center}
\vspace{-0.4em}
\begin{center}{\small\itshape Figure~7.1. Les cinq phases de la mission et
leurs livrables}\end{center}
\end{minipage}

\medskip

\begin{tableaucours}{L{2.6cm}L{4.8cm}L{4.8cm}}{%
    \tableauHeader{Phase} &
    \tableauHeader{Activités principales} &
    \tableauHeader{Chapitre}
}
0. Préparation & Mandat, termes de référence, documents sollicités,
équipe & Chapitre~\ref{chap:planifier} \\
1. Planification & Prise de connaissance, cartographie préliminaire des
risques, décomposition de l'espace d'audit, plan d'audit &
Chapitres~\ref{chap:planifier} et~\ref{chap:decomposition} \\
2. Exécution & Entretiens, tests, cotation \OUI{}, \NON{} ou \NEK{},
papiers de travail & Chapitre~\ref{chap:preuve} \\
3. Rapport & Calcul de $\mu$, $U$, $H$, observations, contradictoire,
rapport & Chapitre~\ref{chap:conclure} \\
4. Suivi & Suivi des recommandations, préparation du cycle suivant &
Chapitre~\ref{chap:suivi} \\
\end{tableaucours}
\vspace{-0.8em}
\begin{center}{\small\itshape Tableau~7.4. Les phases de la mission}\end{center}

Le cycle ne se referme pas sur lui-même~: le suivi d'une mission prépare la
suivante. C'est la dimension temporelle de l'audit, que l'auditologie pose
en axiome~: le système audité évolue, et son état à la mission suivante
dépend des interventions déclenchées par la précédente (axiome~5,
\emph{individuation temporelle}~\cite{minka2026auditologie}).

\subsection{Entrées et sorties de chaque phase}

\begin{tableaucours}{L{2.4cm}L{4.6cm}L{5.2cm}}{%
    \tableauHeader{Phase} &
    \tableauHeader{Elle ne commence pas sans} &
    \tableauHeader{Elle ne se termine pas sans}
}
0 & Un commanditaire identifié & Lettre de mission signée, liste O3 remise \\
1 & La lettre de mission & Plan d'audit validé, décomposition revue \\
2 & Le plan d'audit & Toutes les fiches d'assertion cotées, papiers de
travail revus \\
3 & Les papiers de travail revus & Réponses des audités consignées, rapport
signé \\
4 & Le rapport et ses recommandations & État du suivi présenté au comité
d'audit \\
\end{tableaucours}
\vspace{-0.8em}
\begin{center}{\small\itshape Tableau~7.5. Les conditions d'entrée et de
sortie des phases}\end{center}

\subsection{La mission d'AUDIX CONSEIL chez LiZbiZ}

\begin{boxscenario}{Opération LIZBIZ. Le calendrier de la seconde mission}
\begin{tableaucours}{L{3.0cm}L{2.6cm}L{6.2cm}}{%
    \tableauHeader{Phase} &
    \tableauHeader{Dates (2026)} &
    \tableauHeader{Faits marquants}
}
0. Préparation & 4 au 8 mai & Lettre de mission signée~; 21 documents
sollicités \\
1. Planification & 11 au 22 mai & Prise de connaissance des quatre centres
de données~; décomposition du domaine D~; plan d'audit validé par le
directeur général \\
2. Exécution & 25 mai au 26 juin & Entretiens, extractions, tests sur
quatre agences \\
3. Rapport & 29 juin au 17 juillet & Réunion de restitution, réponses
écrites des audités, rapport final \\
4. Suivi & À partir de janvier~2027 & Revue de la mise en œuvre des
recommandations \\
\end{tableaucours}

\medskip
\textbf{Analyse} (confiance élevée). Le calendrier réserve trois semaines
au rapport et au contradictoire~: c'est le temps nécessaire pour que les
conclusions soient discutées, et donc opposables.
\end{boxscenario}

\subsection{Le budget de la mission}

Le chef de mission traduit le calendrier en jours d'auditeur. Pour la
seconde mission de LiZbiZ~: 5~jours de préparation, 20 de planification,
75 d'exécution (quatre sites), 20 de rapport et de contradictoire, soit
120~jours-auditeur. L'exécution concentre environ 60~\% du budget~; si
elle déborde, ce n'est jamais le contradictoire qu'on raccourcit, mais le
périmètre qu'on réduit, par écrit, avec l'accord du commanditaire.

% ============================================================================
\section{La Boîte à Outils de la Méthode}
\label{sec:c07-outils}
% ============================================================================

La méthode du cours est outillée~: à chaque phase correspondent un ou
plusieurs instruments, numérotés de O0 à O9. Ils seront construits au fil
des chapitres et rassemblés en annexe.

\begin{tableaucours}{cL{3.6cm}L{2.4cm}L{5.4cm}}{%
    \tableauHeader{Outil} &
    \tableauHeader{Nom} &
    \tableauHeader{Phase} &
    \tableauHeader{Usage}
}
O0 & Aide-mémoire du cycle & Toutes & Entrées, activités et livrables de
chaque phase \\
O1 & Fiche technique de la mission & 0 & Présentation de la mission~:
objectifs, acteurs, référentiels \\
O2 & Termes de référence & 0 & Objectifs, périmètre, équipe, livrables,
calendrier \\
O3 & Liste des documents sollicités & 0 & Documents attendus, priorité,
interlocuteur~; un document essentiel absent est coté \NEK{} \\
O4 & Classeur d'aide à l'audit & 1 à 3 & Décomposition selon le gabarit
canonique, fiches d'assertion, cotation \OUI{}, \NON{} ou \NEK{}, calcul de
$\mu$, $U$ et $H$ par domaine \\
O5 & Guide d'entretien & 2 & Questions issues des fiches d'assertion,
classées par interlocuteur \\
O6 & Plan d'audit & 1 & Contexte, périmètre, risques, programme de
travail, calendrier \\
O7 & Plan de tests & 2 & Tests codifiés, preuves attendues, références des
papiers de travail \\
O8 & Projets d'observation & 3 & Observations en sept éléments, soumises
au contradictoire \\
O9 & Suivi des documents & 2 & Mise à jour, sur le terrain, de la liste O3 \\
\end{tableaucours}
\vspace{-0.8em}
\begin{center}{\small\itshape Tableau~7.6. Les outils O0 à O9}\end{center}

L'outil central est le classeur~O4~: il porte la décomposition de l'espace
d'audit (Chapitre~\ref{chap:decomposition}), la cotation des réponses
(Chapitre~\ref{chap:preuve}) et les trois mesures de chaque domaine
(Chapitre~\ref{chap:conclure}).

% ============================================================================
\section{Les Papiers de Travail}
\label{sec:c07-papiers}
% ============================================================================

\begin{boxdef}
\termecle{Papiers de travail}\index{papiers de travail}~: ensemble des
documents établis ou obtenus par l'auditeur pendant la mission, qui
décrivent les travaux réalisés, les preuves recueillies et les conclusions
tirées. Ils permettent à un auditeur expérimenté, extérieur à la mission, de
comprendre ce qui a été fait et d'aboutir aux mêmes conclusions.
\end{boxdef}

\begin{boxprocedure}
\textbf{Ce que contient un papier de travail}
\begin{itemize}
    \item Une référence unique (par exemple PT-21), et le renvoi à la fiche
          d'assertion testée~;
    \item l'objectif du test, le critère, la population et l'échantillon~;
    \item la description du test, assez précise pour être refait~;
    \item les preuves, ou leur référence (fichier, empreinte SHA-256,
          lettre d'extraction)~;
    \item la constatation, séparée de la conclusion, et la cote~;
    \item le nom de l'auteur, la date, le nom du réviseur et la date de
          revue.
\end{itemize}
\end{boxprocedure}

\begin{boxinfo}
\textbf{Le critère du tiers.} Le Chapitre~\ref{chap:prologue} l'a montré~:
le commissaire aux comptes a demandé les papiers de travail de la première
mission, et il n'en existait pas. Un dossier se juge par une question~: un
tiers qui n'a pas participé à la mission peut-il refaire le raisonnement
et arriver à la même cote~? C'est la forme opérationnelle de l'axiome
d'opposabilité.
\end{boxinfo}

\begin{boxmetier}{Auditeur SI}
\textbf{Ouvrir une mission}
\begin{enumerate}
    \item Vérifier que la charte ou la lettre de mission donne l'accès aux
          personnes, aux documents et aux systèmes nécessaires.
    \item Tenir une réunion d'ouverture avec le commanditaire et les
          principaux audités~: objectifs, périmètre, calendrier, règles du
          contradictoire.
    \item Annoncer que la réponse \NEK{} est admise et documentée, et qu'un
          document non remis sera coté comme tel.
    \item Remettre la liste des documents sollicités (O3) avec une date de
          remise.
    \item Consigner la réunion dans un compte rendu, premier papier de
          travail de la mission.
\end{enumerate}
\end{boxmetier}

% ============================================================================
\section{Le Contradictoire}
\label{sec:c07-contradictoire}
% ============================================================================

\subsection{Pourquoi le contradictoire~?}

\begin{boxdef}
\termecle{Contradictoire}\index{contradictoire}~: principe selon lequel
l'audité prend connaissance des constats qui le concernent et peut y
répondre, avant que le rapport ne soit définitif~; ses réponses sont
consignées et, le cas échéant, reproduites dans le rapport.
\end{boxdef}

Le contradictoire n'est pas une politesse. Il remplit trois fonctions~:
\begin{enumerate}
    \item \textbf{Corriger les erreurs} de l'auditeur, qui a pu mal
          comprendre un processus ou manquer une pièce.
    \item \textbf{Rendre les conclusions opposables}~: un constat discuté et
          maintenu est plus difficile à contester qu'un constat découvert
          dans le rapport (axiome~6).
    \item \textbf{Réduire l'incertitude}~: l'audité apporte souvent, au
          moment du contradictoire, l'information qui manquait.
\end{enumerate}

\subsection{Le contradictoire réduit l'incertitude}

\begin{boxlecture}{Domaine B, avant et après la restitution}
Douze questions du domaine~B (gestion des sauvegardes) sont cotées avant la
réunion de restitution, puis après, lorsque le responsable informatique
remet les journaux de sauvegarde qu'il n'avait pas fournis.

\medskip
\begin{tableaucours}{L{3.8cm}ccccc}{%
    \tableauHeader{Moment} &
    \tableauHeader{\OUI{}} & \tableauHeader{\NON{}} & \tableauHeader{\NEK{}} &
    \tableauHeader{$U$} & \tableauHeader{$H$ (bit)}
}
Avant la restitution & 6 & 2 & 4 & 0,33 & 1,46 \\
Après la restitution & 8 & 3 & 1 & 0,08 & 1,19 \\
\end{tableaucours}

\medskip
\textbf{Constatation.} Le taux d'incertitude passe de 0,33 à 0,08~: trois
\NEK{} sont levés, deux en \OUI{}, un en \NON{}.

\medskip
\textbf{Analyse} (confiance élevée). Le contradictoire a surtout réduit
l'ignorance. L'entropie baisse moins, parce que les réponses restent
partagées entre \OUI{} et \NON{}~: le domaine est désormais mieux connu,
et l'on sait où il est défaillant. Sans restitution, le rapport aurait
conclu sur un tiers de réponses inconnues.

\medskip
\hfill\lecture{0{,}67}{0{,}08}{1{,}19}
\end{boxlecture}

\begin{boxprocedure}
\textbf{La trace du contradictoire.} Pour qu'un constat soit opposable, le
dossier de la mission contient~:
\begin{itemize}
    \item le projet d'observation transmis à l'audité, daté~;
    \item la réponse écrite de l'audité, ou le compte rendu de la réunion de
          restitution signé des participants~;
    \item la position finale de l'auditeur~: constat maintenu, modifié ou
          retiré, et pourquoi~;
    \item les pièces nouvelles remises au moment du contradictoire, avec leur
          date de remise.
\end{itemize}
\end{boxprocedure}

\begin{boxalert}
\textbf{Violation typique.} Accepter, au contradictoire, une affirmation de
l'audité («~les sauvegardes sont testées chaque mois~») et changer un
\NON{} en \OUI{} sans pièce.

\smallskip
\textbf{Règle~:} au contradictoire, une affirmation sans pièce peut lever un
malentendu, jamais remplacer une preuve. Elle transforme au mieux un \NON{}
en \NEK{}, en attendant la pièce.
\end{boxalert}

\subsection{Quand l'audité refuse}

Le contradictoire n'oblige pas l'auditeur à l'accord. Si l'audité conteste
un constat sans apporter de pièce, l'auditeur maintient le constat et
reproduit la position de l'audité dans le rapport, à côté de la sienne. Si
l'audité refuse de répondre dans le délai, le rapport le mentionne~: le
silence n'est ni une approbation ni un démenti.

\begin{boxmetier}{RSSI}
\textbf{Se préparer à être audité}
\begin{enumerate}
    \item Remettre les documents sollicités dans les délais~: un document
          non remis sera coté \NEK{}, et signalé comme tel.
    \item Désigner un interlocuteur unique pour l'équipe d'audit.
    \item Répondre par écrit aux projets d'observation, pièce à l'appui.
    \item Proposer, pour chaque recommandation acceptée, un responsable et
          une échéance.
\end{enumerate}
\end{boxmetier}

\begin{boxmetier}{Contrôleur interne}
\textbf{Être l'interlocuteur de l'audit externe.} Chez LiZbiZ,
Mme~NKENG coordonne la remise des documents, assiste aux réunions de
restitution et tient le tableau des recommandations. Elle n'est pas
auditée sur ce rôle, mais son indépendance vis-à-vis des directions
opérationnelles conditionne la qualité des pièces remises.
\end{boxmetier}

\begin{boxinfo}
\textbf{Ce que dit la recherche.} L'axiome d'opposabilité d'A0 dit qu'une
conclusion n'a de valeur que si elle peut être défendue devant un tiers.
Le contradictoire en est l'épreuve anticipée~: il expose la conclusion à la
partie qui a le plus intérêt à la contester, avant qu'un juge, un
régulateur ou un comité ne le fasse~\cite{minka2026auditologie}.
\end{boxinfo}

\separateur

\begin{boxresume}
\textbf{Ce qu'il faut retenir de ce chapitre~:}
\begin{itemize}
    \item \textbf{Sept principes} guident l'auditeur selon l'ISO~19011:2018,
          dont l'approche fondée sur la preuve et l'approche par les risques.
    \item \textbf{La charte d'audit} fonde les pouvoirs et l'indépendance de
          l'audit~; elle est approuvée par l'organe de gouvernance. La
          lettre de mission dit quoi, selon quoi, comment et pour qui.
    \item \textbf{L'équipe} répartit les rôles~; chaque papier de travail est
          revu par un plus expérimenté~; le chef de mission répond des
          travaux de l'expert.
    \item \textbf{Cinq phases}~: préparation, planification, exécution,
          rapport, suivi~; chacune a ses conditions d'entrée et de sortie.
    \item \textbf{Dix outils} (O0 à O9) outillent le cycle~; le classeur O4
          porte la décomposition, la cotation et les mesures.
    \item \textbf{Les papiers de travail} permettent à un tiers de refaire
          le raisonnement.
    \item \textbf{Le contradictoire} corrige, rend opposable et réduit
          l'incertitude~; une affirmation sans pièce ne remplace jamais une
          preuve.
\end{itemize}
\end{boxresume}

\begin{boxinfo}
\textbf{Points de contrôle de l'élève.} Je sais~: citer les sept
principes~; rédiger les six rubriques d'une charte~; dire qui revoit quoi
dans une équipe~; donner l'entrée et la sortie de chaque phase~; lister le
contenu d'un papier de travail~; tracer un contradictoire.
\end{boxinfo}

\bigskip

\begin{boxexo}[title={\ding{45}\quad Exercice~7.1. Quel principe~? \niveaubase}]
Indiquez le principe de l'ISO~19011 que chacune des situations suivantes
méconnaît.
\begin{enumerate}[(a)]
    \item Un auditeur évoque, lors d'un dîner, les failles du système de
          paiement de la BANQUE SANAGA.
    \item Un rapport passe sous silence un constat qui met en cause le
          directeur général.
    \item Un auditeur consacre autant de temps aux imprimantes qu'aux
          comptes à privilèges.
    \item Un constat repose sur ce qu'un agent a « entendu dire ».
    \item Un consultant audite la politique de sécurité qu'il a rédigée
          l'année précédente.
\end{enumerate}
\end{boxexo}

\begin{boxexo}[title={\ding{45}\quad Exercice~7.2. La charte d'audit de LiZbiZ \niveaumoyen}]
LiZbiZ veut créer une fonction d'audit interne.
\begin{enumerate}[(a)]
    \item Rédigez les six rubriques de sa charte d'audit.
    \item À qui la directrice de l'audit doit-elle être rattachée, et
          pourquoi~?
    \item Quels pouvoirs d'accès faut-il prévoir pour auditer l'application
          intégrée et les quatre centres de données~?
\end{enumerate}
\end{boxexo}

\begin{boxexo}[title={\ding{45}\quad Exercice~7.3. Entrées et sorties \niveaubase}]
Pour chacune des situations, dites quelle condition du Tableau~7.5 n'est
pas remplie~: (a) l'équipe commence les entretiens avant la validation du
plan~; (b) le rapport est envoyé avant la réponse des audités~; (c) une
fiche d'assertion n'est pas cotée à la fin du terrain.
\end{boxexo}

\begin{boxexo}[title={\ding{45}\quad Exercice~7.4. Revoir un papier de travail \niveaumoyen}]
Un auditeur junior remet le papier de travail suivant~: « PT-14. Revue des
comptes. Vu le responsable informatique, qui affirme que les comptes sont
revus chaque trimestre. Conclusion~: conforme. » Rédigez les remarques du
réviseur et le papier de travail corrigé, en précisant la cote que vous
pouvez attribuer en l'état.
\end{boxexo}

\begin{boxexo}[title={\ding{45}\quad Exercice~7.5. Le budget \niveaumoyen}]
La mission de LiZbiZ prend du retard~: au 15~juin, 70 des 75~jours
d'exécution sont consommés et le site de Bafoussam n'est pas audité. Le
directeur général propose de réduire le contradictoire à une semaine.
Que répondez-vous~? Proposez deux solutions et la forme de l'accord à
obtenir.
\end{boxexo}

\begin{boxexo}[title={\ding{45}\quad Exercice~7.6. Le contradictoire chiffré \niveaumoyen}]
Recalculez $\mu$, $U$ et $H$ du domaine~B avant et après la restitution
(posez les calculs). Expliquez pourquoi $H$ baisse moins que $U$.
\end{boxexo}

\begin{boxexo}[title={\ding{45}\quad Exercice~7.7. L'expert \niveaumoyen}]
L'ingénieur réseau d'AUDIX CONSEIL remet un rapport de 40~pages issu d'un
outil de scan, qui conclut « 312 vulnérabilités ». Que fait le chef de
mission avant d'utiliser ce rapport~? Comment le résultat entre-t-il dans
les fiches d'assertion du domaine~D~?
\end{boxexo}

\begin{boxexo}[title={\ding{45}\quad Exercice~7.8. L'audité qui refuse \niveaumoyen}]
Le responsable de l'agence de Garoua conteste par écrit un constat sur le
tableur parallèle de facturation, sans pièce. Rédigez le paragraphe du
rapport qui présente le constat et la position de l'audité.
\end{boxexo}

\begin{boxexo}[title={\ding{45}\quad Exercice~7.9. Le contradictoire à la BANQUE SANAGA \niveauexpert}]
Mme~EKAMBI présente à la DSI un projet d'observation~: aucun test de retour
arrière n'a été réalisé avant la bascule de mars~2026. La DSI répond que des
tests ont eu lieu, mais que leurs comptes rendus ont été perdus.
\begin{enumerate}[(a)]
    \item Quelle cotation retenir pour ce contrôle, avant et après la
          réponse de la DSI~? Justifiez.
    \item Que doit contenir le dossier pour que le constat final soit
          opposable~?
    \item Rédigez la position finale de l'auditrice.
    \item En quoi cette situation illustre-t-elle la différence entre réduire
          l'incertitude et lever un constat~?
\end{enumerate}
\end{boxexo}

\begin{boxexo}[title={\ding{45}\quad Exercice~7.10. Concilier deux principes \niveauexpert}]
Le rapport d'audit de la BANQUE SANAGA décrit une faille permettant de
modifier un virement après sa validation. Le rapport doit être transmis au
conseil, à la COBAC et, en résumé, aux agences. Comment conciliez-vous la
confidentialité et la présentation impartiale~? Proposez le contenu de
chaque version.
\end{boxexo}

\bigskip

\begin{boxinfo}
\textbf{Pour aller plus loin.}
\begin{itemize}
    \item ISO~19011:2018~\cite{iso19011_2018}~: les principes, la gestion
          d'un programme d'audit et la compétence des auditeurs.
    \item IIA, Normes mondiales d'audit interne~\cite{iia2024gias}~: le
          domaine III sur la gouvernance de la fonction d'audit.
    \item ISACA, ITAF~\cite{isaca2020itaf}~: les exigences de
          documentation et de supervision propres à l'audit des SI.
\end{itemize}
\end{boxinfo}

\begin{boxinfo}
\textbf{Transition.} La mission commence avant le terrain. Le
Chapitre~\ref{chap:planifier} détaille la préparation et la planification~:
termes de référence, documents sollicités, prise de connaissance du système
d'information et plan d'audit.
\end{boxinfo}
